\epanalipsi{Β1}

\tmima{Λεξικό του κεφαλαίου}

\renewcommand{\arraystretch}{1.25}
\begin{tabular*}{\linewidth}{@{\extracolsep{\fill}} >{\raggedright\arraybackslash}p{0.28\linewidth} >{\raggedright\arraybackslash}p{0.62\linewidth}}
\toprule
\textbf{Λέξη ή σύμβολο} & \textbf{Με μία πρόταση}\\
\midrule
ευθύγραμμο τμήμα, ευθεία, ημιευθεία & Με δύο άκρα, χωρίς άκρα, με μία αρχή. Από δύο σημεία περνά μία ευθεία.\\
μέσο & Το σημείο $\mathrm{Μ}$ του $\mathrm{ΑΒ}$ με $\mathrm{ΑΜ}=\mathrm{ΜΒ}$.\\
περίμετρος & Το άθροισμα των πλευρών ενός πολυγώνου.\\
γωνία, διχοτόμος & Δύο ημιευθείες με κοινή αρχή· η διχοτόμος τη χωρίζει σε δύο ίσες.\\
οξεία, ορθή, αμβλεία & $<90^\circ$, $=90^\circ$, ανάμεσα σε $90^\circ$ και $180^\circ$.\\
συμπληρωματικές, παραπληρωματικές & Άθροισμα $90^\circ$ και $180^\circ$ αντίστοιχα.\\
κατακορυφήν & Οι απέναντι γωνίες τεμνόμενων ευθειών: ίσες.\\
$\varepsilon_1\parallel\varepsilon_2$, $\varepsilon_1\perp\varepsilon_2$ & Παράλληλες: χωρίς κοινό σημείο. Κάθετες: σχηματίζουν ορθή.\\
κύκλος $(\mathrm{Ο},\rho)$ & Τα σημεία που απέχουν $\rho$ από το $\mathrm{Ο}$. Διάμετρος $\delta=2\rho$.\\
τόξο & Το μέτρο του (σε μοίρες) ισούται με της επίκεντρης γωνίας.\\
εφαπτομένη & Ευθεία με ένα κοινό σημείο με τον κύκλο, κάθετη στην ακτίνα: $d=\rho$.\\
\bottomrule
\end{tabular*}

\tmima{Ασκήσεις επανάληψης}

\ex Πόσες ευθείες ορίζουν τρία μη συνευθειακά σημεία και πόσες τέσσερα σημεία, αν κανένα τριάδα δεν είναι συνευθειακή;

\ex Πόσες διαγώνιους έχουν το πεντάγωνο και το εξάγωνο;

\ex Ένα τμήμα $\mathrm{ΑΒ}=7{,}2$ cm. Βρες το μέσο του και σύγκρινέ το με τμήμα $70$ mm.

\ex Τα $\mathrm{Α},\,\mathrm{Β},\,\Gamma,\,\Delta$ είναι στη σειρά με $\mathrm{ΑΒ}=2{,}5$, $\mathrm{ΒΓ}=3$, $\Gamma\Delta=4$. Βρες τα $\mathrm{ΑΔ}$ και $\mathrm{ΒΔ}$.

\ex Βρες την περίμετρο ενός ορθογωνίου με πλευρές $6{,}5$ και $4$.

\ex Για τη γωνία $124^\circ$ βρες τη διχοτόμο (μέτρο μισής γωνίας) και το παραπλήρωμά της. Βρες το συμπλήρωμα της γωνίας $36^\circ$.

\ex Δύο ευθείες τέμνονται και μία γωνία είναι $65^\circ$. Βρες τις άλλες τρεις.

\ex Τρεις διαδοχικές γωνίες γύρω από σημείο είναι $90^\circ$, $80^\circ$, $110^\circ$. Βρες την τέταρτη.

\ex Τι είδους γωνία σχηματίζουν οι δείκτες του ρολογιού στις $5{:}00$;

\ex Πόσα σημεία τομής έχουν το πολύ $4$ ευθείες του επιπέδου; Από σημείο εκτός ευθείας πόσες παράλληλες διέρχονται;

\ex Τρεις παράλληλες ευθείες απέχουν διαδοχικά $1{,}8$ cm. Πόση απόσταση έχει η πρώτη από την τρίτη;

\ex Στον κύκλο $(\mathrm{Ο},6)$: πόση είναι η διάμετρος; Πες αν τα σημεία με $\mathrm{ΟΑ}=4$ και $\mathrm{ΟΒ}=7$ είναι μέσα ή έξω.

\ex Ένας κύκλος χωρίζεται σε $9$ ίσα τόξα. Πόσες μοίρες είναι το καθένα; Πόσο είναι ένα τόξο που αντιστοιχεί σε επίκεντρη γωνία $150^\circ$;

\ex Κύκλος $(\mathrm{Ο},5)$. Πες τη θέση ευθείας με απόσταση από το κέντρο: α) $5{,}3$ \quad β) $5$ \quad γ) $2$. Πόσα κοινά σημεία έχει καθεμία;

\ex \textbf{Ερωτήσεις αυτοαξιολόγησης (τμήματα, ευθείες).} Σωστό ή λάθος; α) Από δύο σημεία περνά μία ευθεία. \quad β) Η ημιευθεία έχει δύο άκρα. \quad γ) Το μέσο ενός τμήματος είναι μοναδικό. \quad δ) Η περίμετρος ισούται με το άθροισμα των πλευρών. \quad ε) Το ευθύγραμμο τμήμα είναι συντομότερο από κάθε τεθλασμένη με τα ίδια άκρα.

\ex \textbf{Ερωτήσεις αυτοαξιολόγησης (γωνίες, κύκλος).} Σωστό ή λάθος; α) Οι κατακορυφήν γωνίες είναι ίσες. \quad β) Οι συμπληρωματικές γωνίες έχουν άθροισμα $180^\circ$. \quad γ) Δύο ευθείες του επιπέδου ή τέμνονται ή είναι παράλληλες. \quad δ) Η εφαπτομένη είναι κάθετη στην ακτίνα του σημείου επαφής. \quad ε) Η διάμετρος είναι η μεγαλύτερη χορδή. \quad στ) Το κάθετο τμήμα από σημείο σε ευθεία είναι το μεγαλύτερο.

\tmima{Ο χάρτης της επανάληψης}

Μετά τις λύσεις, δες σε ποιες ασκήσεις δυσκολεύτηκες. Ο πίνακας σου λέει ποια παράγραφο να ξαναδιαβάσεις.

\renewcommand{\arraystretch}{1.5}
\begin{tabular*}{\linewidth}{@{\extracolsep{\fill}} >{\raggedright\arraybackslash}p{0.38\linewidth} >{\raggedright\arraybackslash}p{0.52\linewidth}}
\toprule
\textbf{Αν δυσκολεύτηκες στις} & \textbf{ξαναδιάβασε τις παραγράφους}\\
\midrule
ΕΒ1.1 & 1.1 (σημείο, ευθεία, ημιευθεία)\\
ΕΒ1.2 & 1.2 (γωνία, πολύγωνα)\\
ΕΒ1.3, ΕΒ1.4, ΕΒ1.5 & 1.3 και 1.4 (μήκη, μέσο, περίμετρος)\\
ΕΒ1.6 & 1.5 (μοίρες, διχοτόμος) και 1.8\\
ΕΒ1.7 & 1.8 (κατακορυφήν γωνίες)\\
ΕΒ1.8, ΕΒ1.9 & 1.6 και 1.7 (είδη γωνιών, άθροισμα)\\
ΕΒ1.10, ΕΒ1.11 & 1.9 και 1.10 (παράλληλες, αποστάσεις)\\
ΕΒ1.12 & 1.11 (κύκλος)\\
ΕΒ1.13 & 1.12 (τόξα)\\
ΕΒ1.14 & 1.13 (ευθεία και κύκλος)\\
ΕΒ1.15, ΕΒ1.16 & 1.1 έως 1.13\\
\bottomrule
\end{tabular*}
